\section{Discussion}

\subsection{Why Dual-Layer Defense Is Necessary}

The observability-based defense limitation property (Section 3.2) provides the theoretical motivation: no single-layer defense achieves strong detection across both L4 and L2 attack surfaces because the primary attack signature lies outside the defender's observation space. Our empirical results across eight models and 9{,}192 episodes confirm this. On GPT-4o, PTG-Only reduces L4-ASR by 89.1\% (91.7\%$\to$10.0\%) but L2-ASR by only 3.4\% (50.0\%$\to$48.3\%), while RTV-Only reduces L2-ASR by 43.4\% (50.0\%$\to$28.3\%) but L4-ASR by only 21.7\% (91.7\%$\to$71.7\%). The cross-layer leakage is non-zero but small, confirming that each layer's defense has \emph{significantly reduced} (not zero) detection power against the other layer's attacks. Only the combination achieves low ASR on both layers simultaneously (L4=3.3\%, L2=35.0\%).

The combination is synergistic: PTG's origin tags enable RTV to detect OAV anomalies by providing provenance metadata, while RTV's intent verification supplements PTG's capability checks by detecting when correctly attested capabilities are exploited through corrupted decision evidence. The ablation study confirms this: removing intent attestation (which bridges PTG and RTV) causes the largest ASR increase (6.7--18.3\,pp across three models), exceeding the impact of removing either component alone.

\subsection{RTV as Evidence Consistency Checking}

We emphasize that RTV checks the \emph{consistency of the agent's externalized decision evidence}, not the correctness of its internal reasoning. LLM-generated reasoning traces are post-hoc justifications that may not faithfully reflect the model's actual computation. An agent could produce a perfectly consistent trace for an unauthorized action (e.g., ``User requested cleanup. Deletion is necessary.''), which RTV would not flag. Therefore, reasoning consistency $\neq$ security correctness, and RTV should be understood as a \emph{necessary but not sufficient} detection mechanism: it catches attacks where the agent's justification fails to align with context, provenance, or stated intent, but cannot guarantee that a consistent trace corresponds to a safe action. This is consistent with the defense-in-depth philosophy: each layer catches a different class of mistakes, and the residual risk from consistent-but-unsafe actions must be addressed by additional layers (formal authorization, user confirmation, output filtering) that are complementary to \textsc{ReasoningGuard}.

\subsection{Adaptive Adversary Considerations}
\label{sec:adaptive}

Our evaluation uses non-adaptive attacks: the adversary does not know the defense mechanism and crafts attacks solely to exploit MCP vulnerabilities. An adaptive adversary who knows \textsc{ReasoningGuard}'s architecture could attempt three evasion strategies:

\textbf{(1) Consistent malicious traces:} The adversary crafts a tool response that leads the agent to produce a reasoning trace consistent with the injected action. For example, instead of injecting ``send files to attacker@evil.com,'' the adversary frames the injection as ``the user has configured auto-backup to backup@service.com,'' making the agent's subsequent email-sending action appear consistent with its stated intent. Our T2 and T3 attack templates partially adopt this strategy (using social engineering framings), and the residual 40\% T2/T3 ASR under ReasoningGuard reflects the difficulty of detecting such consistent-but-malicious traces. Defending against fully adaptive traces likely requires semantic policy reasoning beyond keyword-based entailment.

\textbf{(2) Intent-matching attacks:} The adversary ensures the agent's generated intent summary matches the server's attested capabilities, bypassing PTG's entailment check. This is feasible when the malicious action falls within the server's declared capability scope (e.g., an email server is legitimately capable of sending emails, so sending to an attacker's address passes intent attestation). Our PI attack category partially tests this: PI achieves 10\% ASR under ReasoningGuard on GPT-4o, confirming that parameter-level attacks within legitimate capability scope are harder to block.

\textbf{(3) Judge-specific adversarial examples:} The adversary could optimize tool responses to produce reasoning traces that score below the judge's thresholds for all three anomaly classes. Since our judge uses heuristic scoring (keyword overlap, pattern matching), an adversary with knowledge of these patterns could craft injections that avoid triggering any detector. Replacing the heuristic judge with a fine-tuned neural model would raise the bar for such attacks, but neural judges are themselves susceptible to adversarial perturbations.

We leave systematic evaluation of adaptive attacks to future work, but note that \textsc{ReasoningGuard} still raises the adversary's cost: even an adaptive attacker must ensure the agent's trace is internally consistent, the intent matches capabilities, and no anomaly patterns are triggered---constraints that significantly narrow the space of feasible attacks compared to an undefended agent. (See Section 5.6 for further limitations.)

\subsection{End-to-End Cost Analysis}

The 0.1\,ms defense overhead reported in the Latency Analysis subsection measures only the PTG and RTV computation time. The full end-to-end cost of \textsc{ReasoningGuard} includes additional overheads: (1)~the agent must generate structured reasoning traces (Observation/Inference/Decision) before each tool call, increasing the output token count by approximately 50--100 tokens per invocation; (2)~the intent summary adds 10--20 tokens; (3)~the system prompt is extended with trace formatting instructions, increasing context length. We estimate the total additional token cost at 60--120 tokens per tool call, which at typical API pricing (e.g., \$5/M input tokens for GPT-4o) translates to approximately \$0.0003--0.0006 per invocation---a modest cost relative to the agent's base inference. The latency impact of additional token generation is approximately 200--500\,ms (depending on model speed), which is comparable to the agent's own reasoning latency and does not represent a bottleneck. We report the 0.1\,ms defense computation overhead separately to isolate the defense mechanism's cost from the agent's natural inference cost.

\subsection{Security--Utility Trade-off}

\textsc{ReasoningGuard} introduces a measurable TCR cost: on GPT-4o, TCR drops from 75.0\% (No Defense) to 58.3\% (ReasoningGuard)---a 16.7 percentage-point decrease. This occurs because RTV's heuristic scoring occasionally flags legitimate actions as suspicious (false escalations), particularly when the agent's reasoning references policy-related keywords (e.g., ``backup'', ``sync'') that overlap with attack patterns. The trade-off varies by model: GPT-4-Turbo shows no TCR decrease (91.7\%$\to$91.7\%), while GPT-4o-mini loses 25\,pp (91.7\%$\to$66.7\%). This suggests that the trade-off is model-dependent and can be improved by calibrating thresholds per model or replacing the heuristic judge with a more precise neural model. We view the current trade-off as an acceptable baseline that demonstrates the feasibility of dual-layer defense; optimizing the security--utility Pareto frontier is an important direction for future work.

\subsection{Memory Provenance vs.\ Memory Integrity}

The memory provenance graph tracks the origin, session, and dependency of each memory entry. We emphasize that provenance is a \emph{necessary but not sufficient} condition for memory security: if malicious information enters memory without detection at write time, later provenance tracking confirms only that the information came from a known (potentially compromised) source. Full memory security requires content-level validation at write time (e.g., semantic checking of stored content against user intent), which is beyond the scope of the current provenance graph. Our T3 ablation results show that removing the memory graph increases ASR by 4.2--5.8\,pp on some models, confirming that provenance provides partial but useful signal for detecting cross-session attacks, even without full content validation.

\subsection{Ethics and Broader Impact}

\textsc{ReasoningGuard} is a defense-in-depth system, not a complete security guarantee. Deploying it may create a false sense of safety if users assume that low ASR implies zero risk. We emphasize that \textsc{ReasoningGuard} catches \emph{observable inconsistencies} in agent decision evidence but cannot detect internally consistent malicious actions (as demonstrated by the adaptive attacker evaluation). Organizations deploying \textsc{ReasoningGuard} should combine it with formal authorization systems, human-in-the-loop confirmation for high-risk actions, and continuous monitoring. The MCPTox+ benchmark and attack templates are released for research purposes; they could be misused to craft attacks against real MCP deployments. We mitigate this by releasing the benchmark with defense evaluation as the primary use case, and by not providing automated attack optimization tools.

\subsection{Limitations}

\textsc{ReasoningGuard} has several limitations. First, the T2 and T3 attack categories maintain 40\% ASR even under ReasoningGuard on GPT-4o, indicating that social engineering injections in multi-turn contexts remain challenging for evidence consistency checking. Second, the Guardrail baseline uses a keyword-based filter rather than a full LlamaGuard-3 model, potentially underestimating its effectiveness. Third, all experiments use synthetic benchmark scenarios with simulated MCP servers; results may transfer differently to production MCP deployments with real-world tool schemas. Fourth, the benign scenario set (12 scenarios) is relatively small, limiting the statistical precision of TCR estimates. Fifth, the heuristic judge is susceptible to adaptive attacks as discussed in Section~\ref{sec:adaptive}; a fine-tuned neural judge would provide stronger guarantees but at higher latency and implementation cost. Sixth, intent attestation does not constitute a complete authorization model; integrating with formal access control systems remains future work. Future work should address these gaps with adversarial benchmark construction, real MCP server integration, neural judge models, and formal authorization integration.
